\documentclass[11pt]{article}
\usepackage{mathblatt}
% Lösungen nach mathe-nachhilfe gemeinsam.md „Lösungen“ und pruefung.md § 6.
\newcommand{\kennung}{NUB}
\begin{document}
\pfheftstil
\pfheftkopf[\kennung]{Symmetrie}{P10 \textperiodcentered{} Lösungen}

\begin{pfloesung}{}
\lz{1a)}{$1$ Achse}{Höhe auf die Grundseite}{}
\lz{b)}{$2$ Achsen}{die beiden Diagonalen}{}
\lz{c)}{$3$ Achsen}{je von einer Ecke zur Mitte der Gegenseite}{}
\lz{d)}{$0$ Achsen}{keine zwei Seiten gleich lang}{}
\end{pfloesung}
\begin{pfloesung}{}
\lz{2.}{B}{Parallelogramm: nur punktsymmetrisch; Rechteck 2, Raute 2 Achsen}{}
\end{pfloesung}
\begin{pfloesung}{}
\lz{3.}{B}{$1$ Achse: senkrecht durch die Mitten der parallelen Seiten}{}
\end{pfloesung}
\begin{pfloesung}{}
\lz{4.}{B}{$1$ Achse: die Diagonale $\overline{BD}$; sie halbiert $\overline{AC}$, nicht umgekehrt}{}
\end{pfloesung}
\end{document}
